\section*{Chapter 2 Exercises}

These exercises fill in the calculations we deferred in the main text. None of
them require any tools beyond what's in the chapter.

% ---------------------------------------------------------------------
\subsection*{2.1 The loss in covariance form.}
In the setup, we claimed that ``with some matrix algebra'' the mean squared error
can be rewritten purely in terms of the data covariances. Verify this. Starting
from
\begin{align*}
    \Lcal(W_1, \dots, W_L)
    = \mathbb{E}_{(x,y)\sim\Dcal}
      \Big[\Tr\,(y - W_L \cdots W_1 x)(y - W_L \cdots W_1 x)^\top \Big],
\end{align*}
expand the product inside the trace, use linearity of the trace and the
expectation, and substitute the definitions $\Sigma_{xx} := \mathbb{E}_\Dcal[xx^\top]$,
$\Sigma_{yx} := \mathbb{E}_\Dcal[yx^\top]$, and $\Sigma_{yy} := \mathbb{E}_\Dcal[yy^\top]$
to arrive at
\begin{align*}
    \Lcal(W_1, \dots, W_L)
    = \|\Sigma_{yx}\Sigma_{xx}^{-1/2} - W_L \cdots W_1 \Sigma_{xx}^{1/2}\|_\mathrm{F}^2
      + \mathrm{const.}
\end{align*}
where $\mathrm{const.} = \Sigma_{yy} - \Sigma_{yx}(\Sigma_{xx})^{-1}(\Sigma_{yx})^\top$.
Confirm that the constant term does not depend on the weight matrices.

\emph{Hint:} ``complete the square'' at the level of matrices---add and subtract
$\Sigma_{yx}\Sigma_{xx}^{-1/2}$ inside the Frobenius norm.

% ---------------------------------------------------------------------
% \begin{exercise}[The conservation law, by the algebra trick.]
% Derive the conservation law of Equation~(\ref{eq:conservation-ex}) the same way the
% chapter body does.
% \begin{enumerate}[label=(\alph*)]
%     \item Compute the gradients $\nabla_{W_1}\Lcal$ and $\nabla_{W_2}\Lcal$ for the
%     2-layer loss and confirm the gradient flow equations
%     \begin{align*}
%         \frac{d}{dt}W_1 &= W_2^\top(\Sigma_{yx} - W_2 W_1 \Sigma_{xx}), &
%         \frac{d}{dt}W_2 &= (\Sigma_{yx} - W_2 W_1 \Sigma_{xx}) W_1^\top.
%     \end{align*}
%     \item Using these, show that
%     $\big(\tfrac{d}{dt}W_1\big)W_1^\top = W_2^\top\big(\tfrac{d}{dt}W_2\big)$, and
%     hence that
%     \begin{equation}
%         \frac{d}{dt}\big(W_2^\top W_2 - W_1 W_1^\top\big) = 0.
%         \label{eq:conservation-ex}
%     \end{equation}
%     \item Conclude that $H := W_2^\top W_2 - W_1 W_1^\top$ is constant along the
%     gradient flow trajectory, so a network initialized balanced ($H \approx 0$)
%     stays balanced.
% \end{enumerate}
% \end{exercise}

% ---------------------------------------------------------------------
\subsection*{2.2 The conservation law, from the symmetry.}
The algebra trick used in the chapter to derive the conservation law $\frac{d}{dt}[W_2^\top W_2 - W_1 W_1^\top] = 0$ derives it as if by
luck. This exercise re-derives it from the continuous
symmetry of the loss, making clear \emph{why} the trick works.

Throughout, use the matrix inner product $\langle A, B\rangle := \Tr(A^\top B)$,
and recall the gradient flow $\dot\theta = -\nabla\Lcal$ for $\theta = (W_1, W_2)$.

\begin{enumerate}[label=(\alph*)]
    \item \emph{(Symmetry $\Rightarrow$ orthogonality.)} Suppose $\Lcal$ is
    invariant under a one-parameter family of transformations
    $\theta \mapsto \theta + \epsilon\, v(\theta)$, with generator (velocity field)
    $v$. Show that invariance, to first order in $\epsilon$, is equivalent to
    \begin{align*}
        \langle v(\theta), \nabla\Lcal(\theta)\rangle = 0 \quad \text{for all } \theta.
    \end{align*}
    \item \emph{(Gradient flow never moves along the symmetry.)} Substitute the
    gradient flow equation to conclude that $\langle v(\theta), \dot\theta\rangle = 0$
    along any trajectory. Interpret this geometrically: the gradient flow velocity
    is always orthogonal to the symmetry direction, so optimization never spends
    motion traveling along the symmetry orbit.
    \item \emph{(The specific symmetry.)} For the deep linear network, the symmetry
    is $W_1 \mapsto AW_1,\ W_2 \mapsto W_2 A^{-1}$. Take $A = e^{\epsilon M} \approx
    I + \epsilon M$ for an arbitrary generator matrix $M$, and show that the
    induced generator is $v = (M W_1,\ -W_2 M)$.
    \item \emph{(All generators at once.)} Plug this $v$ into the orthogonality
    condition from (b), substitute $\nabla_{W_1}\Lcal = -\dot W_1$ and
    $\nabla_{W_2}\Lcal = -\dot W_2$, and use the cyclic property of the trace to
    rewrite both terms with $M^\top$ pulled to the front. Show the condition
    becomes
    \begin{align*}
        \Tr\!\big(M^\top \dot W_1 W_1^\top\big) = \Tr\!\big(M^\top W_2^\top \dot W_2\big).
    \end{align*}
    \item \emph{(Conclude.)} Since this must hold for \emph{every} generator $M$,
    deduce that $\dot W_1 W_1^\top = W_2^\top \dot W_2$, which is exactly
    $\frac{d}{dt}(W_2^\top W_2 - W_1 W_1^\top) = 0$. Note where the arbitrariness of
    $M$ was essential: it is what upgrades a single scalar orthogonality condition
    into the full matrix-valued conservation law.
\end{enumerate}

\emph{Remark.} Compare this derivation to the algebra trick. The trick yields
$\dot W_1 W_1^\top = W_2^\top \dot W_2$ by a fortunate cancellation; here the
\emph{same} equation appears as the direct consequence of ``the flow is orthogonal
to every symmetry direction.'' The luck of the trick was the symmetry in disguise
all along.

% ---------------------------------------------------------------------
\subsection*{2.3 Alignment collapses the model to its magnitudes.}
Write the weight SVDs as $W_1 = U_1 S_1 V_1^\top$ and $W_2 = U_2 S_2 V_2^\top$, and
recall the rotated model $\tilde{W}_2 \tilde{W}_1 = U^{\star\top} W_2 W_1 V^\star$,
where $\Sigma_{yx} = U^\star \Lambda^\star {V^\star}^\top$. Show that under the
three weight-alignment conditions,
\begin{align*}
    U_2 = U^\star, \qquad V_1 = V^\star, \qquad V_2 = U_1,
\end{align*}
the rotated model reduces to the product of magnitude matrices,
$\tilde{W}_2 \tilde{W}_1 = S_2 S_1$, with all rotation factors canceling. Identify
which orthogonality relation kills each rotation factor.

% ---------------------------------------------------------------------
\subsection*{2.4 Solving the scalar learning dynamics.}
This is the climactic integration of the chapter. The $\mu$-th end-to-end singular
value obeys the separable ODE
\begin{align*}
    \frac{d}{dt}s_\mu^2 = 2 s_\mu^2 (\lambda^\star_\mu - s_\mu^2).
\end{align*}
Let $u := s_\mu^2$ for brevity, so that $\dot{u} = 2u(\lambda^\star_\mu - u)$.
\begin{enumerate}[label=(\alph*)]
    \item Separate variables and use the partial-fraction identity
    \begin{align*}
        \frac{1}{u(\lambda^\star_\mu - u)}
        = \frac{1}{\lambda^\star_\mu}\left(\frac{1}{u} + \frac{1}{\lambda^\star_\mu - u}\right)
    \end{align*}
    to integrate both sides from $0$ to $t$.
    \item Solve the resulting expression for $u = s_\mu^2(t)$ and confirm you
    recover the logistic solution
    \begin{align*}
        s^2_\mu(t) = \frac{\lambda^\star_\mu}
        {1 + \left(\dfrac{\lambda^\star_\mu}{s^2_\mu(0)} - 1\right) e^{-2 \lambda^\star_\mu t}}.
    \end{align*}
    \item Check the two limits: confirm that $s_\mu^2(t) \to s_\mu^2(0)$ as
    $t \to 0$ and $s_\mu^2(t) \to \lambda^\star_\mu$ as $t \to \infty$.
\end{enumerate}

% ---------------------------------------------------------------------
\subsection*{2.5 Time-to-learn.}
Using the logistic solution from the previous exercise, derive the halfway-time
$\tau_\mu$ at which mode $\mu$ reaches half its final strength,
$s^2_\mu(\tau_\mu) = \tfrac{1}{2}\lambda^\star_\mu$. Show that, dropping order-one
constants,
\begin{align*}
    \tau_\mu \approx \frac{1}{\lambda^\star_\mu} \ln \frac{\lambda^\star_\mu}{s^2_\mu(0)}.
\end{align*}
Then argue from this expression that (a) larger modes are learned first, and (b)
the time-separation between consecutive modes widens as the initialization scale
$s^2_\mu(0)$ shrinks---the mechanism behind the cleanly separated stairsteps in
the loss curve.
